% pline-doc.tex
% arara: lualatex: { options: ["-shell-escape"] }
% arara: lualatex: { options: ["-shell-escape"] }
\PassOptionsToPackage{english}{babel}
\PassOptionsToPackage{colorlinks = true}{hyperref}
\PassOptionsToPackage{margin = 1in}{geometry}
\PassOptionsToPackage{depth = 4}{bookmark}
\documentclass[12pt]{report}
\usepackage{
    pline
    ,microtype
    ,babel
    ,fontspec
    ,unicode-math,
    ,caption
    ,subcaption
    ,xcolor
    ,geometry
    ,minted
    ,bookmark
    ,hyperref
}
\setmathfont{IBMPlexMath-Regular}
\counterwithin{figure}{section}
\counterwithin{table}{section}
\DeclareCaptionFormat{pline}{#1#2\par\noindent#3}
\captionsetup{
    justification = raggedright
    ,singlelinecheck = false
    ,format = pline
}
\captionsetup[subfigure]{format=hang}
\setcounter{tocdepth}{3}
\setmainfont{Carlito}
\title{
    The pline Package,\\
    Version 0.01\\
    {\small https://github.com/Pseudonym321/TikZ-Animations/tree/master1/TikZ/pline}
}
\author{Jasper}
\date{\today}
\begin{document}
    \maketitle 
    \tableofcontents
    \chapter{Introduction}
        pline is a Lua-based 3D graphics package for 
        Ti\textit{k}Z.
        \section{What are the current capabilities of pline?}
            pline can render multiple non-intersecting 
            triangulated mesh surfaces at once. It is also possible 
            to draw a plane which has been clipped by a cube.
            There are commands for drawing and labelling geometric 
            vectors too. Curves and surfaces can be defined using 
            matrix transformations.
        \section{
            How does pline improve on existing packages?
        }
            There are two significant packages for 3D artwork in 
            Ti\textit{k}Z: pgfplots and tikz-3dplot.
            pgfplots can z-sort the faces of a single parametric 
            surface, but not multiple surfaces at once.
            Both pgfplots and tikz-3dplot are capable of changing 
            the camera view, but the result is necessarily 
            orthographic, and the rotations are restricted to 
            azimuth and elevation; that is, the \(z\)-axis is always 
            vertical. Neither pgfplots nor tikz-3dplot enable 
            matrix  pline is built to render 
            multiple parametric objects at once, enable arbitrary
            object orientations in non-orthographic projections, 
            and enable matrix 
        \section{Where is the package at, and where it is heading?}
            The vision of the package pline is to one 
            day be able to render arbitrary numbers of intersecting 
            polygons, by clipping them with each other; this would 
            enable arbitrary surface intersections, without 
            visual artefacts. A more near term goal of 
            pline is to also implement an ODE 
            solving tool which approximates trajectories in 
            vector fields. I also want to add capabilities for 
            doing programming in Lua, instead of in pgf.
            I also want to add culling for projective 
            transformations.
    \chapter{Geometric Vectors}
        \section{How are geometric vectors defined in pline?}
            Geometric vectors are defined by two things: a 
            start-coordinate, and a relative displacement.
            There are two classifications of geometric vectors;
            there are zero vectors, and non-zero vectors.
            A zero vector is represented by a small circle.
            A non-zero vector is represented by a line segment
            with an arrow. The endpoints of a non-zerovector 
            may or may not have a circle centered on the endpoint.
            If there is no circle drawn, the arrow stops directly
            at the endpoint. If a circle is drawn, then the arrow 
            tip is moved back by the circle's radius. Observe 
            the tips of the vectors in Figure \ref{fig:2-1-1}.
            Both tips lie on the same point, despite one arrow tip 
            being offset by the circle's radius.
            \begin{figure}
                \centering
                \begin{tikzpicture}
                    \drawvector[points = both,green]{0,0}{2,0}
                    \drawvector[blue]{4,0}{-2,0}
                \end{tikzpicture}
                \caption{The exact location of a geometric vector's tip.}
                \label{fig:2-1-1}
            \end{figure}
            The syntax for drawing Figure \ref{fig:2-1-1} is shown in 
            the following code.
\begin{minted}{tex}
\begin{tikzpicture}
    \drawvector[points = both,yellow]{0,0}{2,0}
    \drawvector[blue]{4,0}{-2,0}
\end{tikzpicture}
\end{minted}
            There are four values for the key 
            \mintinline{latex}{points}: \mintinline{latex}{neither}, 
            \mintinline{latex}{behind}, \mintinline{latex}{front} and 
            \mintinline{latex}{both}. These tell the command
            \mintinline{latex}{\drawvector} whether and where to put 
            circles over the endpoints. See Figure  \ref{fig:2-1-2}, 
            which uses  the following code.
\begin{minted}{tex}
% a)
\begin{tikzpicture}
    \drawvector[points = neither]{0,0}{0,1}
\end{tikzpicture}
% b)
\begin{tikzpicture}
    \drawvector[points = behind]{1,0}{0,1}
\end{tikzpicture}
% c)
\begin{tikzpicture}
    \drawvector[points = front]{2,0}{0,1}
\end{tikzpicture}
% d)
\begin{tikzpicture}
    \drawvector[points = both]{3,0}{0,1}
\end{tikzpicture}
\end{minted}
            \begin{figure}
                \centering
                \begin{subfigure}{0.3\textwidth}
                    \centering
                    \begin{tikzpicture}
                        \drawvector[points = neither]{0,0}{0,1}
                    \end{tikzpicture}
                    \caption{\mintinline{latex}{points = neither}}
                \end{subfigure}
                \begin{subfigure}{0.3\textwidth}
                    \centering
                    \begin{tikzpicture}
                        \drawvector[points = behind]{1,0}{0,1}
                    \end{tikzpicture}
                    \caption{\mintinline{latex}{points = behind}}
                \end{subfigure}\par
                \begin{subfigure}{0.3\textwidth}
                    \centering
                    \begin{tikzpicture}
                        \drawvector[points = front]{2,0}{0,1}
                    \end{tikzpicture}
                    \caption{\mintinline{latex}{points = front}}
                \end{subfigure}
                \begin{subfigure}{0.3\textwidth}
                    \centering
                    \begin{tikzpicture}
                        \drawvector[points = both]{3,0}{0,1}
                    \end{tikzpicture}
                    \caption{\mintinline{latex}{points = both}}
                \end{subfigure}
                \caption{The values for the key \mintinline{latex}{points}.}
                \label{fig:2-1-2}
            \end{figure}
            The \mintinline{latex}{\drawvector} syntax is mostly
            meant for diagrams which don't use \(z\)-sorting.
            Still, they can be used in some 3D still-diagrams;
            for instance, these vectors connect nicely.
            They will also eventually be able to be drawn flat 
            on planes. There is also a zero-vector command 
            available via \mintinline{latex}{\drawpoint}.
            See Figure \ref{fig:2-1-3}, with the following code.
\begin{minted}{tex}
\begin{tikzpicture}
    \drawpoint[red]{0,0}
\end{tikzpicture}
\end{minted}
            \begin{figure}
                \centering
                \begin{tikzpicture}
                    \drawpoint[red]{0,0}
                \end{tikzpicture}
                \caption{Zero-vector}
                \label{fig:2-1-3}
            \end{figure}
    \chapter{Parametric Objects}
        \section{How is a parametric object defined?}
            Parametric objects are defined using a handful of 
            different pieces of information. In particular, they 
            use parametric functions, domain parameters and a samples
            parameter for each dimension in the domain. Figure 
            \ref{fig:3-1-1} is defined by the following code:
\begin{minted}{tex}
\begin{tikzpicture}
    \appendcurve[
        u min = 0
        ,u max = {tau}
        ,u samples = 200
        ,transformation = {euler(pi/2,pi/3,pi/2)}
        ,x = {2*cos(u)}
        ,y = {2*sin(u)}
        ,z = {u}
        ,draw options = {
            purple
            ,ultra thick
            ,line join = round
        }
    ]
    \appendsurface[
        u min = 0
        ,u max = {tau}
        ,u samples = 36
        ,v min = 0
        ,v max = {tau}
        ,v samples = 36*2/3
        ,transformation = {euler(pi/2,pi/3,pi/2)}
        ,x = {2*cos(u)+sphere(u,v)[1][1]}
        ,y = {2*sin(u)+sphere(u,v)[1][2]}
        ,z = {u+sphere(u,v)[1][3]}
        ,draw options = {
            draw
            ,blue,ultra thin
            ,line join = round
        }
        ,fill options = {
            fill = green
            ,fill opacity = 0.7
        }
    ]
    \foreach \u in {0,...,35} {
        \foreach \v in {0,...,35} {
            \pgfmathsetmacro{\uscale}{2*pi/(36-1)}
            \pgfmathsetmacro{\vscale}{2*pi/(36*2/3-1)}
            \pgfmathsetmacro{\u}{\u*\uscale}
            \pgfmathsetmacro{\v}{\v*\vscale}
            \appendcurve[
                u min = 0
                ,u max = 0.3
                ,u samples = 2
                ,transformation = {euler(pi/2,pi/3,pi/2)}
                ,x = {
                    2 * 
                    cos(token.get_macro("u")) +
                    (u + 1) * 
                    sphere(
                        token.get_macro("u")
                        ,token.get_macro("v")
                    )[1][1]
                }
                ,y = {
                    2 * 
                    sin(token.get_macro("u")) +
                    (u + 1) * 
                    sphere(
                        token.get_macro("u")
                        ,token.get_macro("v")
                    )[1][2]
                }
                ,z = {
                    token.get_macro("u") + 
                    (u + 1) * 
                    sphere(
                        token.get_macro("u")
                        ,token.get_macro("v")
                    )[1][3]
                }
                ,draw options = {-{Stealth[round]}}
            ]
        } 
    }
    \rendersegments
\end{tikzpicture}
\end{minted}
            \begin{figure}
                \centering
                \begin{tikzpicture}
                    \appendcurve[
                        u min = 0
                        ,u max = {tau}
                        ,u samples = 200
                        ,transformation = {euler(pi/2,pi/3,pi/2)}
                        ,x = {2*cos(u)}
                        ,y = {2*sin(u)}
                        ,z = {u}
                        ,draw options = {
                            purple
                            ,ultra thick
                            ,line join = round
                        }
                    ]
                    \appendsurface[
                        u min = 0
                        ,u max = {tau}
                        ,u samples = 36
                        ,v min = 0
                        ,v max = {tau}
                        ,v samples = 36*2/3
                        ,transformation = {euler(pi/2,pi/3,pi/2)}
                        ,x = {2*cos(u)+sphere(u,v)[1][1]}
                        ,y = {2*sin(u)+sphere(u,v)[1][2]}
                        ,z = {u+sphere(u,v)[1][3]}
                        ,draw options = {
                            draw
                            ,blue,ultra thin
                            ,line join = round
                        }
                        ,fill options = {
                            fill = green
                            ,fill opacity = 0.7
                        }
                    ]
                    \foreach \u in {0,...,35} {
                        \foreach \v in {0,...,35} {
                            \pgfmathsetmacro{\uscale}{2*pi/(36-1)}
                            \pgfmathsetmacro{\vscale}{2*pi/(36*2/3-1)}
                            \pgfmathsetmacro{\u}{\u*\uscale}
                            \pgfmathsetmacro{\v}{\v*\vscale}
                            \appendcurve[
                                u min = 0
                                ,u max = 0.3
                                ,u samples = 2
                                ,transformation = 
                                    {euler(pi/2,pi/3,pi/2)}
                                ,x = {
                                    2 * 
                                    cos(token.get_macro("u")) +
                                    (u + 1) * 
                                    sphere(
                                        token.get_macro("u")
                                        ,token.get_macro("v")
                                    )[1][1]
                                }
                                ,y = {
                                    2 * 
                                    sin(token.get_macro("u")) +
                                    (u + 1) * 
                                    sphere(
                                        token.get_macro("u")
                                        ,token.get_macro("v")
                                    )[1][2]
                                }
                                ,z = {
                                    token.get_macro("u") + 
                                    (u + 1) * 
                                    sphere(
                                        token.get_macro("u")
                                        ,token.get_macro("v")
                                    )[1][3]
                                }
                                ,draw options = {-{Stealth[round]}}
                            ]
                        } 
                    }
                    \rendersegments
                \end{tikzpicture}
                \caption{A triangulated mesh surface and a curve.}
                \label{fig:3-1-1}
            \end{figure}
            The package is capable of handling division by zero errors
            without erroring.
            The code for Figure \ref{fig:3-1-2} is as follows.
\begin{minted}{tex}
\begin{tikzpicture}
    \clip (-5,-5) rectangle (5,5);
    \pgfmathsetmacro{\angle}{0}
    \let\angle\angle
    \foreach \longitude in {1,...,36}{
        \let\longitude\longitude
        \appendcurve[
            u min = 0
            ,u max = {pi}
            ,u samples = 4*45
            ,transformation = {euler(pi/2,pi/3,pi/6)}
            ,x = {
                stereographic_projection(
                    matrix_multiply(
                        sphere(
                            tau * 
                            token.get_macro("longitude") /
                            36
                            ,u
                        )
                        ,euler(
                            0
                            ,pi/2
                            ,(2*pi/24) *
                            token.get_macro("angle") / 
                            36
                        )
                    )
                )[1][1]
            }
            ,y = {
                stereographic_projection(
                    matrix_multiply(
                        sphere(
                            tau *
                            token.get_macro("longitude") /
                            36
                            ,u
                        )
                        ,euler(
                            0
                            ,pi/2
                            ,(2*pi/24) * 
                            token.get_macro("angle") / 
                            36
                        )
                    )
                )[1][2]
            }
            ,z = 0*u
        ]
        \appendcurve[
            u min = 0
            ,u max = {pi}
            ,u samples = 23
            ,transformation = {
                matrix_multiply(
                    euler(
                        0
                        ,pi/2
                        ,(2*pi/24) * 
                        token.get_macro("angle") / 
                        36
                    )
                    ,euler(pi/2,pi/3,pi/6)
                )
            }
            ,x = {
                sphere(
                    tau * 
                    token.get_macro("longitude") / 
                    36
                    ,u
                )[1][1]
            }
            ,y = {
                sphere(
                    tau * 
                    token.get_macro("longitude") / 
                    36
                    ,u
                )[1][2]
            }
            ,z = {
                sphere(
                    tau * 
                    token.get_macro("longitude") / 
                    36
                    ,u
                )[1][3]
            }
        ]
    }
    \foreach \latitude in {1,...,18}{
        \let\latitude\latitude
        \appendcurve[
            u min = 0
            ,u max = {2*pi}
            ,u samples = 4*45
            ,transformation = {euler(pi/2,pi/3,pi/6)}
            ,x = {
                stereographic_projection(
                    matrix_multiply(
                        sphere(
                            u
                            ,tau *
                            token.get_macro("latitude") / 
                            36
                        )
                        ,euler(
                            0
                            ,pi/2
                            ,(tau/24) * 
                            token.get_macro("angle") / 
                            36
                        )
                    )
                )[1][1]
            }
            ,y = {
                stereographic_projection(
                    matrix_multiply(
                        sphere(
                            u
                            ,tau *
                            token.get_macro("latitude") /
                            36
                        )
                        ,euler(
                            0
                            ,pi/2
                            ,(tau/24) *
                            token.get_macro("angle") /
                            36
                        )
                    )
                )[1][2]
            }
            ,z = 0*u
        ]
        \appendcurve[
            u min = 0
            ,u max = {2*pi}
            ,u samples = 45
            ,transformation = {
                matrix_multiply(
                    euler(
                        0
                        ,pi/2
                        ,(2*pi/24) *
                        token.get_macro("angle") /
                        36
                    )
                    ,euler(pi/2,pi/3,pi/6)
                )
            }
            ,x = {
                sphere(
                    u
                    ,tau*token.get_macro("latitude")/36
                )[1][1]
            }
            ,y = {
                sphere(
                    u
                    ,tau*token.get_macro("latitude")/36
                )[1][2]
            }
            ,z = {
                sphere(
                    u
                    ,tau*token.get_macro("latitude")/36
                )[1][3]
            }
        ]
    }
    \rendersegments
\end{tikzpicture}
\end{minted}
            \begin{figure}
                \centering
                \begin{tikzpicture}
                    \clip (-5,-5) rectangle (5,5);
                    \pgfmathsetmacro{\angle}{0}
                    \let\angle\angle
                    \foreach \longitude in {1,...,36}{
                        \let\longitude\longitude
                        \appendcurve[
                            u min = 0
                            ,u max = {pi}
                            ,u samples = 4*45
                            ,transformation = {euler(pi/2,pi/3,pi/6)}
                            ,x = {
                                stereographic_projection(
                                    matrix_multiply(
                                        sphere(
                                            tau * 
                                            token.get_macro("longitude") /
                                            36
                                            ,u
                                        )
                                        ,euler(
                                            0
                                            ,pi/2
                                            ,(2*pi/24) *
                                            token.get_macro("angle") / 
                                            36
                                        )
                                    )
                                )[1][1]
                            }
                            ,y = {
                                stereographic_projection(
                                    matrix_multiply(
                                        sphere(
                                            tau *
                                            token.get_macro("longitude") /
                                            36
                                            ,u
                                        )
                                        ,euler(
                                            0
                                            ,pi/2
                                            ,(2*pi/24) * 
                                            token.get_macro("angle") / 
                                            36
                                        )
                                    )
                                )[1][2]
                            }
                            ,z = 0*u
                        ]
                        \appendcurve[
                            u min = 0
                            ,u max = {pi}
                            ,u samples = 23
                            ,transformation = {
                                matrix_multiply(
                                    euler(
                                        0
                                        ,pi/2
                                        ,(2*pi/24) * 
                                        token.get_macro("angle") / 
                                        36
                                    )
                                    ,euler(pi/2,pi/3,pi/6)
                                )
                            }
                            ,x = {
                                sphere(
                                    tau * 
                                    token.get_macro("longitude") / 
                                    36
                                    ,u
                                )[1][1]
                            }
                            ,y = {
                                sphere(
                                    tau * 
                                    token.get_macro("longitude") / 
                                    36
                                    ,u
                                )[1][2]
                            }
                            ,z = {
                                sphere(
                                    tau * 
                                    token.get_macro("longitude") / 
                                    36
                                    ,u
                                )[1][3]
                            }
                        ]
                    }
                    \foreach \latitude in {1,...,18}{
                        \let\latitude\latitude
                        \appendcurve[
                            u min = 0
                            ,u max = {2*pi}
                            ,u samples = 4*45
                            ,transformation = {euler(pi/2,pi/3,pi/6)}
                            ,x = {
                                stereographic_projection(
                                    matrix_multiply(
                                        sphere(
                                            u
                                            ,tau *
                                            token.get_macro("latitude") / 
                                            36
                                        )
                                        ,euler(
                                            0
                                            ,pi/2
                                            ,(tau/24) * 
                                            token.get_macro("angle") / 
                                            36
                                        )
                                    )
                                )[1][1]
                            }
                            ,y = {
                                stereographic_projection(
                                    matrix_multiply(
                                        sphere(
                                            u
                                            ,tau *
                                            token.get_macro("latitude") /
                                            36
                                        )
                                        ,euler(
                                            0
                                            ,pi/2
                                            ,(tau/24) *
                                            token.get_macro("angle") /
                                            36
                                        )
                                    )
                                )[1][2]
                            }
                            ,z = 0*u
                        ]
                        \appendcurve[
                            u min = 0
                            ,u max = {2*pi}
                            ,u samples = 45
                            ,transformation = {
                                matrix_multiply(
                                    euler(
                                        0
                                        ,pi/2
                                        ,(2*pi/24) *
                                        token.get_macro("angle") /
                                        36
                                    )
                                    ,euler(pi/2,pi/3,pi/6)
                                )
                            }
                            ,x = {
                                sphere(
                                    u
                                    ,tau*token.get_macro("latitude")/36
                                )[1][1]
                            }
                            ,y = {
                                sphere(
                                    u
                                    ,tau*token.get_macro("latitude")/36
                                )[1][2]
                            }
                            ,z = {
                                sphere(
                                    u
                                    ,tau*token.get_macro("latitude")/36
                                )[1][3]
                            }
                        ]
                    }
                    \rendersegments
                \end{tikzpicture}
                \caption{Automatic division by zero handling.}
                \label{fig:3-1-2}
            \end{figure}
            \begin{figure}
                \centering
                \begin{tikzpicture}
                    \clip (-5,-5) rectangle (5,5);
                    \foreach[count = \c from 1] \angle in {20,40,...,80} {
                        \pgfmathsetmacro{\angle}{\angle*pi/180}
                        \let\angle\angle
                        \appendsurface[
                            u min = pi/6
                            ,u max = 0.75*tau
                            ,u samples = 36 
                            ,v min = 0
                            ,v max = tau
                            ,v samples = 36
                            ,transformation = {euler(pi/2,pi/3,pi/6+pi)}
                            ,draw options = {draw = black}
                            ,fill options = {
                                \ifnum\c=1
                                    fill = red
                                \fi
                                \ifnum\c=2
                                    fill = yellow
                                \fi
                                \ifnum\c=3
                                    fill = green
                                \fi
                                \ifnum\c=4
                                    fill = blue
                                \fi
                                ,fill opacity = 0.6
                            }
                            ,x = {
                                (1 / cos(token.get_macro("angle"))) *
                                cos(u) + 
                                sqrt(
                                    (
                                        1 / 
                                        cos(token.get_macro("angle"))
                                    )^2 - cos(token.get_macro("angle"))
                                ) * 
                                sphere(u,v)[1][1]
                            }
                            ,y = {
                                (1 / cos(token.get_macro("angle"))) *
                                sin(u) + 
                                sqrt(
                                    (
                                        1 / 
                                        cos(token.get_macro("angle"))
                                    )^2 - cos(token.get_macro("angle"))
                                ) * 
                                sphere(u,v)[1][2]
                            }
                            ,z = {
                                sqrt(
                                    (
                                        1 / 
                                        cos(token.get_macro("angle"))
                                    )^2 - cos(token.get_macro("angle"))
                                ) * 
                                sphere(u,v)[1][3]
                            }
                        ]
                    }
                    \rendersegments
                \end{tikzpicture}
                \caption{Stereographically nested tori.}
                \label{fig:3-1-3}
            \end{figure}
    \chapter{Clipped Subspaces}
        \section{Clipping individual planes.}
            \begin{figure}
                \centering
                \begin{tikzpicture}
                    \appendprism[%
                        a = 1
                        ,b = 2
                        ,c = 3
                        ,d = 4
                        ,x min = -1
                        ,x max = 1
                        ,y min = -1
                        ,y max = 1
                        ,z min = -1
                        ,z max = 1
                        ,transformation = {euler(pi/2,pi/3,pi/3)}
                    ]
                    \appendplane[%
                        a = 1
                        ,b = 2
                        ,c = 3
                        ,d = 1
                        ,x min = -1
                        ,x max = 1
                        ,y min = -1
                        ,y max = 1
                        ,z min = -1
                        ,z max = 1
                        ,transformation = {euler(pi/2,pi/3,pi/3)}
                    ]
                    \rendersegments
                \end{tikzpicture}
                \caption{%
                    A plane which has been clipped by a rectangular 
                    prism.
                }
                \label{fig:4-1-1}
            \end{figure}
\end{document}